\section{KẾ HOẠCH GIAI ĐOẠN CUỐI VÀ KẾT LUẬN}\label{sec:ke_hoach}

\subsection{Kế hoạch công việc còn lại}

Bảng~\ref{tab:ke_hoach_cuoi} liệt kê các công việc từ mốc giữa kỳ đến ngày bảo vệ dự kiến (cuối tháng 11/2026).

\begin{table}[H]
    \centering
    \small
    \caption{Kế hoạch giai đoạn cuối}
    \label{tab:ke_hoach_cuoi}
    \begin{tabularx}{\textwidth}{|c|X|P{3.0cm}|P{3.6cm}|}
        \hline
        \textbf{\#} & \textbf{Công việc} & \textbf{Thời gian} & \textbf{Kết quả cần đạt} \\
        \hline
        1 & Đo lặp SD-WAN bằng công cụ web: failover link WAN (Site 100 và chi nhánh, 3--5 lần mỗi kịch bản), thời gian khôi phục khi mất vSmart, độ sẵn sàng dài hạn & 25/10--05/11 & Bảng RTO min/trung bình/max, tỉ lệ gói đến đích, MTTR \\
        \hline
        2 & Kiểm thử failover firewall ASAv HA và VRRP trên Core & 25/10--31/10 & Biên bản kiểm thử có điều kiện, số gói mất và thời gian chuyển \\
        \hline
        3 & Hoàn thiện Site 500: cấu hình vEdge thành Site 500, nối switch và máy trạm, thêm site vào chính sách vSmart & 26/10--05/11 & vEdge-S500 lên fabric; máy trạm nhận DHCP và ping được các site \\
        \hline
        4 & SDN: giảm khựng controller, cải thiện trường hợp hỏng kép và mất liên kết Distribution--Core & 26/10--05/11 & Số đo lặp lại ổn định hơn; ghi rõ giới hạn còn lại \\
        \hline
        5 & Thu bằng chứng nếu toàn bộ vEdge mất kết nối lần nữa & Theo dõi liên tục & Kết luận nguyên nhân hoặc ghi nhận không tái diễn \\
        \hline
        6 & Hoàn thiện báo cáo cuối kỳ & 01/11--18/11 & Báo cáo đầy đủ 5 chương, cập nhật số liệu kiểm thử \\
        \hline
        7 & Slide và kịch bản demo trên lab & 08/11--20/11 & Slide; kịch bản demo đã chạy thử \\
        \hline
    \end{tabularx}
\end{table}

\subsection{Rủi ro và biện pháp}

\begin{itemize}
    \item \textbf{Tài nguyên máy chủ EVE-NG:} lab 72 node trên máy chủ dùng chung; chạy đồng thời dễ làm chậm controller và vEdge. Biện pháp: chỉ bật các site cần cho từng phép thử, tắt cụm SD-WAN khi đo SDN và ngược lại, ghi lại tải máy chủ cùng mỗi kết quả đo.
    \item \textbf{Mất trạng thái khi wipe node:} vEdge mất chứng thư, Windows/OVS mất cấu hình tay. Biện pháp: không wipe các node này; lưu cấu hình và quy trình phục hồi trong kho mã nguồn.
    \item \textbf{Whitelist controller mất sau khởi động lại:} vEdge rơi khỏi fabric mà không có cảnh báo rõ ràng. Biện pháp: kiểm tra whitelist sau mỗi lần khởi động lại controller và đẩy lại từ vManage; cảnh báo email báo ngay khi vEdge down.
    \item \textbf{Thời gian:} khối lượng đo kiểm và viết báo cáo dồn vào tháng 11. Biện pháp: thực hiện kiểm thử song song với viết báo cáo, ghi kết quả ngay sau mỗi phép thử. Nếu Site 500 không kịp, dùng vEdge dự phòng đã ký sẵn tại Site 900 để minh họa việc mở chi nhánh mới.
\end{itemize}

\subsection{Kết luận giữa kỳ}

Đến mốc giữa kỳ, nhóm đã hoàn thành phần thiết kế và khối lượng triển khai cốt lõi của đề tài. Mô hình trên EVE-NG đã thể hiện được các mục tiêu chính: Campus chính phân tầng với gateway dự phòng và firewall HA, SDN điều khiển sáu switch OVS bằng Ryu với cơ chế tự phục hồi và fast-failover đã đo bằng cắt link thật, fabric SD-WAN kết nối bốn site qua hai transport với chính sách chọn đường theo SLA, failover chi nhánh, Internet breakout và cảnh báo email, cùng các dịch vụ mạng truy cập được từ chi nhánh.

Nửa còn lại của dự án tập trung vào đo lặp phía SD-WAN, kiểm thử failover firewall, hoàn thiện Site 500 và tài liệu. Nhóm đánh giá dự án đang đúng tiến độ và đủ điều kiện để hoàn thành các mục tiêu đã đề ra trước ngày bảo vệ.
